\chapter{Catalogue, Variants and Marketplace}
\label{ch:catalogue}

\section{Catalogue and product model}

The seeded catalogue has 2,400 products: 30 hand-written (they carry the featured placements) and 2,370 generated deterministically from a seeded pseudo-random generator, so every environment gets the same data. They belong to 43 product types across nine departments, and carry 104 invented brand names. All pictures are flat illustrations generated by a script in the repository. \tagimpl{}

Products are not uniform records. Each \emph{product type} owns a set of \emph{attribute definitions}, each marked as a \emph{variation} attribute (what variants differ by: colour, storage, RAM) or a \emph{spec} attribute (a technical detail: refresh rate, processor), and optionally as a filter facet. A smartphone therefore has storage and RAM; a running shoe has size, colour, cushioning and terrain. A product stores its spec values as typed JSON keyed by attribute; the type supplies the vocabulary and labels. This lets the product page render a technical table and the search page offer the right filters without special-casing categories. \tagimpl{}

\section{The variant system}

A variant is an exact purchasable configuration. It has a SKU, its \emph{selections} (for example colour, storage and RAM), a price, a list price, a stock count and, when colour differs, its own pictures. Of the 8,860 seeded variants, many products have three variation dimensions. Only valid combinations exist: a 1~TB option is not offered with 8~GB RAM if the catalogue does not sell it. \tagimpl{}

On the product page the customer picks values dimension by dimension. A pure function (\code{catalog/variants.ts}, importable from client components) resolves the selection to a variant; when the exact combination does not exist, it moves to the nearest valid variant instead of showing a dead end. Price, stock, SKU and picture follow the resolved variant. \tagtest{} (variant resolution is unit-tested.)

\begin{figure}[H]
\centering
\begin{tikzpicture}[node distance=7mm]
\node[mod] (a) {Product};
\node[mod, right=of a] (b) {Selections\\\scriptsize colour, storage, RAM};
\node[mod, right=of b] (c) {Variant\\\scriptsize resolution};
\node[mod, right=of c] (d) {SKU};
\node[mod, right=of d] (e) {Offer};
\node[mod, right=of e] (f) {Cart line};
\foreach \s/\t in {a/b,b/c,c/d,d/e,e/f} \draw[arr] (\s) -- (\t);
\end{tikzpicture}
\caption{From product page to cart line. The cart stores the variant and the offer, not a product.}
\label{fig:variantflow}
\end{figure}

\paragraph{Why identity must survive.} The cart line stores the variant identifier and the offer identifier; checkout re-reads both; the order snapshots SKU, option label and seller. Two configurations of the same product are therefore two cart lines, and a line is distinguished from another by variant \emph{and} offer. The customer can see which configuration they are buying at every step, and the order still describes what was bought if the catalogue later changes. \tagimpl{} \tagtest{}

\section{Marketplace offers and the buy box}

Every variant has an implicit \emph{first-party} offer (the variant's own price and stock, sold and fulfilled by ``Cartly''). The \code{offers} table adds 608 seller offers, each with its own price, shipping charge, handling time, fulfilment mode (Cartly-fulfilled or seller-fulfilled) and stock. A cart line records which offer it came from, so the same variant from two sellers is two lines, stock is taken from the chosen offer, and a seller line carries its own shipping. \tagimpl{} \tagtest{}

The \emph{buy box} is the offer pre-selected on the page. Our rule is our own and is not Amazon's proprietary algorithm:

\begin{enumerate}
\item Consider only offers with stock; if none, show the first-party offer as unavailable.
\item Compute each offer's \emph{landed price}: price plus shipping. Offers within 2\% of the lowest landed price are contenders.
\item Among contenders prefer Cartly-fulfilled over seller-fulfilled, then the first-party offer, then shorter handling time, then lower landed price, then offer identifier.
\end{enumerate}

The final tie-break on identifier makes the outcome independent of the order in which offers happen to be read. The product page shows the selected seller, fulfilment, price, availability and delivery window, and lists the others under ``Other sellers'' with the same facts so that choosing one is an informed choice rather than noise. The seller's handling time also pushes the delivery window back (Chapter~\ref{ch:commerce}). \tagtest{} (representative tests: tie-breaks, 2\% tolerance, order independence.)

\begin{lstlisting}[caption={The contender rule, simplified from \code{catalog/offers.ts}.},label={lst:buybox}]
const lowest = Math.min(...buyable.map(landed));
const contenders = buyable.filter((o) => landed(o) <= lowest * 1.02);
// rank: Cartly-fulfilled, first-party, handling time, landed price, offer id
\end{lstlisting}
